% Cover letter draft for the CJFR submission. Build: tectonic -o ../manuscript/build cover-letter.tex
\documentclass[11pt,letterpaper]{letter}
\usepackage[T1]{fontenc}
\usepackage{lmodern}
\usepackage[margin=1in]{geometry}
\usepackage{xcolor}
\usepackage[hidelinks]{hyperref}
\newcommand{\pending}[1]{\textcolor{red!70!black}{[\textsc{Pending:} #1]}}
\signature{Gregory Paradis\\on behalf of all authors}
\address{Faculty of Forestry \& Environmental Stewardship\\University of British Columbia\\2424 Main Mall, Office 4645\\Vancouver, BC V6T 1Z4, Canada\\gregory.paradis@ubc.ca \quad +1 604 822 1890}
\begin{document}
\begin{letter}{Editor-in-Chief\\Canadian Journal of Forest Research}
\opening{Dear Editor,}

Please consider our manuscript ``Rolling-horizon machine scheduling for forest harvesting
operations: planning horizon, re-optimization frequency, and solver choice evaluated with state
carry-forward in FHOPS'' by Rosalia Jaffray, Kathleen Coupland, and Gregory Paradis for
publication as a Research Article in the \emph{Canadian Journal of Forest Research}.

Rolling-horizon optimization is the standard way to schedule harvesting machines over a season,
yet planners have little evidence on how long each planning window should be and how often the
plan should be re-optimized. We evaluate both choices in a full factorial experiment on nine
semi-synthetic British Columbia case studies (234 runs), using the open-source Forest Harvesting
Operations Planning System (FHOPS) with two solution methods. Two features distinguish the study.
First, every window starts from the operational state left by the decisions already implemented:
remaining block volume, inventory staged between harvesting phases, and machine positions.
Second, every run is scored on its stitched, implemented schedule, replayed against the full
problem. Almost every configuration delivered within about 4\% of the best plan found for its
scenario. Extending the horizon beyond about four weeks gained little, daily and weekly
re-optimization were indistinguishable, and only a 14-day commitment period carried a small
penalty. The clearest differences were between solvers, not configurations: cold-started
mixed-integer programming reached its time limit in most large windows, whereas simulated
annealing never failed.

The manuscript re-evaluates, with a corrected procedure, a design first reported in the first
author's MASc thesis (University of British Columbia, 2026). The earlier evaluation used a FHOPS
release that did not carry state between windows and scored runs by the last window's
objective. The present results supersede it, and the manuscript says so explicitly. The FHOPS
software itself is described in a separate manuscript under review at \emph{SoftwareX}. That
manuscript does not report any rolling-horizon results, so there is no overlap in results.

All data, code, and run outputs will be openly available (see the Data availability statement).
The authors used generative AI tools (Claude Opus 5.5 and Kimi K3 for experiment code, Claude
Opus 5.5 for editing, and OpenAI GPT-5.x models for FHOPS development), as disclosed in the
Methods and Acknowledgements; the authors verified all outputs. Gregory Paradis leads development of the
open-source FHOPS software used in the study (a non-financial interest); the authors declare no
other competing interests. The manuscript
is not under consideration elsewhere, and all authors have approved the submission.

\pending{Suggested referees (no collaborators): names, affiliations, e-mail addresses.}

\closing{Sincerely,}
\end{letter}
\end{document}
